\section{沿 Mumford 与 Seshadri 的方法证明射影性}\label{src:9.5}
\begin{epigraph}
然而, 这并不完全是人们能够证明的结论, 这也反映了构造的任何表述所固有的一些技术困难.
\epigraphauthor{Luis \protect\'Alvarez-C\'onsul 与 Alastair King \cite[第 3 页]{CK07}}
\end{epigraph}
我们利用 GIT 将参数化光滑曲线 $C$ 上多重稳定向量丛的模空间 $M_{r,d}^{\mathrm{ps}}$ 构造为射影簇, 从而完成定理~\hyperref[thm:B]{B} 的证明. Mumford 最初利用 GIT 将稳定向量丛的模空间 $M_{r,d}^{\mathrm s}$ 构造为拟射影簇 \cite{Mum63}, C. S. Seshadri 则将其扩充为参数化半稳定向量丛的射影紧化 \cite{Ses67}. 我们沿用 Simpson 的构造 \cite{Sim94a}, 它可方便地推广到奇异曲线与高维簇. 其他论述见 \cite{Ses82}, \cite[\S7]{LP97}, \cite[\S4]{HL10}, \cite[\S5.4]{New78} 和 \cite[\S10]{Muk03}.
\subsection{GIT 设置与主要定理}\label{src:9.5.1-subsection}
固定整数 $r>0$ 和 $d$, 设 $P(z)=rz+r(1-g)+d$ 是 $C$ 上秩 $r$, 次数 $d$ 的凝聚层的 Hilbert 多项式. 选取 $p\in C(\Bbbk)$, 对每个整数 $\nu$, 映射 $F\mapsto F\otimes\mathcal O_C(\nu p)$ 给出同构 $\Coh_{r,d}\cong\Coh_{r,d+r\nu}$, 故可假设 $d$ 任意大. 假设 $d>r(2g-1)$, 则每个秩 $r$, 次数 $d$ 的半稳定向量丛 $F$ 均由整体截面生成, 且满足 $\mathrm H^1(C,F)=0$ (引理~\bookref{9.2.7}). 定义向量空间
\[
V:=\Bbbk^{\oplus P(0)},
\]
其维数为半稳定 $F$ 的 $P(0)=\mathrm h^0(C,F)$. $\mathrm H^0(C,F)$ 的一组基给出等同 $\mathrm H^0(C,F)\cong V$, 并在射影 Quot 概形 $\Quot^P(\mathcal O_C\otimes V/C)$ (定理~\bookref{2.1.3}) 中定义商
\[
[\mathcal O_C\otimes V\twoheadrightarrow F]\in\Quot^P(\mathcal O_C\otimes V/C).
\]
此概形带有 $\mathrm{GL}(V)$ 的作用. 此外, 对 $m\gg0$, 态射
\[
\begin{aligned}
\Quot^P(\mathcal O_C\otimes V/C)&\hookrightarrow\Gr\big(P(m),\mathrm H^0(C,\mathcal O_C(m))\otimes V\big),\\
[\mathcal O_C\otimes V\twoheadrightarrow F]&\longmapsto[\mathrm H^0(C,\mathcal O_C(m))\otimes V\twoheadrightarrow\mathrm H^0(C,F(m))]
\end{aligned}
\]
是 $\mathrm{GL}(V)$-等变闭浸入 (定理~\bookref{2.4.5}). 与 Pl\"ucker 嵌入
\[
\Gr\big(P(m),\mathrm H^0(C,\mathcal O_C(m))\otimes V\big)\hookrightarrow\mathbb P\big(\textstyle\bigwedge^{P(0)}(\mathrm H^0(C,\mathcal O_C(m))\otimes V)\big)
\]
复合, 得到 $\Quot^P(\mathcal O_C\otimes V/C)$ 的 $\mathrm{GL}(V)$-等变射影嵌入. 将其上的 $\mathrm{GL}(V)$-线丛 $L_m$ 定义为限制
\[
L_m:=\mathcal O(1)|_{\Quot^P(\mathcal O_C\otimes V/C)}.
\]
本节的主要定理将 $\Quot^P(\mathcal O_C\otimes V/C)$ 关于 $L_m$ 及 $\mathrm{SL}(V)$ 的作用\footnote{该定理对 $\mathrm{GL}(V)$ 的作用不成立.} 的 GIT 半稳定轨迹 (定义~\bookref{8.2.1}) 与半稳定向量丛 (定义~\bookref{9.1.11}) 等同.
\begin{theorem}\label{src:9.5.1}
设 $C$ 是特征零代数闭域 $\Bbbk$ 上亏格 $g\geq2$ 的光滑连通射影曲线. 对每个秩 $r>0$, 存在整数 $D$ 和 $M$, 使得对所有 $d>D$, $m>M$, 对商 $[\mathcal O_C\otimes V\twoheadrightarrow F]\in\Quot^P(\mathcal O_C\otimes V/C)$, 下列条件等价, 其中 $P$ 是秩 $r$, 次数 $d$ 的向量丛的 Hilbert 多项式, $V=\Bbbk^{\oplus P(0)}$:
\begin{enumerate}[(1)]
\item $[\mathcal O_C\otimes V\twoheadrightarrow F]$ 关于 $L_m$ 和 $\mathrm{SL}(V)$ 是 GIT 半稳定的;
\item $F$ 半稳定且 $\mathrm H^0(C,\mathcal O_C\otimes V)\to\mathrm H^0(C,F)$ 为同构.
\end{enumerate}
\end{theorem}
\begin{remark}\label{src:9.5.2}
该定理等价地可表述如下: 对整数 $r>0$ 和 $d$, 存在整数 $N$ 和 $M$, 使得当 $n>N$, $m>n+M$ 时, 对商 $[\mathcal O_C(-n)\otimes V\twoheadrightarrow F]\in\Quot^P(\mathcal O_C(-n)\otimes V/C)$, 其中 $V=\Bbbk^{\oplus P(n)}$, 下列条件等价:
\begin{enumerate}[(1)]
\item $[\mathcal O_C(-n)\otimes V\twoheadrightarrow F]$ 关于 $L_m$ 和 $\mathrm{SL}(V)$ 是 GIT 半稳定的;
\item $F$ 半稳定且 $\mathrm H^0(C,\mathcal O_C\otimes V)\to\mathrm H^0(C,F(n))$ 为同构.
\end{enumerate}
\end{remark}
\begin{remark}\label{src:9.5.3}
定理~\bookref{9.5.1} 的陈述与证明在正特征下也成立. 我们限制到特征零, 是因为仅发展了线性约化群的 GIT, 而未发展一般约化群的 GIT; 见注~\bookref{7.5.4}.
\end{remark}
\begin{corollary}\label{src:9.5.4}
对每个 $r>0$ 和 $d$, $\cBun_{r,d}^{\mathrm{ss}}$ 有射影模空间 $M_{r,d}^{\mathrm{ps}}$.
\end{corollary}
\begin{proof}
只需将已有结论联系起来:
\begin{itemize}
\item 因为 $\Coh_{r,d}\cong\Coh_{r,d+r\nu}$, 可假设 $d\gg0$.
\item 对 $m\gg0$, 该定理将 $\mathbb G_m$-刚性化 $\cBun_{r,d}^{\mathrm{ss}}\!\sslash\mathbb G_m$ 与 $\boldsymbol\mu_n$-刚性化 $[\Quot^P(\mathcal O_C\otimes V/C)^{\mathrm{ss}}/\mathrm{SL}(V)]\!\sslash\boldsymbol\mu_n$ 等同, 其中 GIT 半稳定轨迹是关于 $L_m$ 和 $\mathrm{SL}(V)$ 的作用而言.
\item 若模叠 $\cBun_{r,d}^{\mathrm{ss}}$, $\cBun_{r,d}^{\mathrm{ss}}\!\sslash\mathbb G_m$ 与 $[\Quot^P(\mathcal O_C\otimes V/C)^{\mathrm{ss}}/\mathrm{SL}(V)]$ 中任一个有良模空间, 则三者均有良模空间, 且这些空间同构.
\item 射影 GIT 的主要定理 (\bookref{8.2.6}) 表明 $[\Quot^P(\mathcal O_C\otimes V/C)^{\mathrm{ss}}/\mathrm{SL}(V)]$ 有射影良模空间.
\end{itemize}
\end{proof}
我们强调, GIT 方法将 $M_{r,d}^{\mathrm{ps}}$ 构造为射影簇, 无需依赖模空间作为固有代数空间的存在性 (定理~\bookref{9.3.11}). 特别地, 得到了半稳定性的开性 (\bookref{9.2.6}), Langton 定理 (\bookref{9.3.1}) 及 $\cBun_{r,d}^{\mathrm{ss}}$ 的 $\Theta$-完备性和 S-完备性 (命题~\bookref{9.3.8}) 的独立证明. 另一方面, 半稳定性的有界性 (\bookref{9.2.8}) 及相关有界性论证 (习题~\bookref{9.2.9}) 对 GIT 构造不可或缺.

定理~\bookref{9.5.1} 的证明在 \S\bookref{9.5.3-subsection} 中给出. 尽管证明相对简短, 且是最容易的模空间 GIT 构造之一 (当然也比 $\overline M_g$ 的 GIT 构造容易), GIT 分析仍然细致, 尤其要妥善安排 $D$ 与 $M$ 的界. 半稳定性判据 (\bookref{9.5.5}) 选取 $D$, 并用对 $E\subseteq F$ 的不等式 $\mathrm h^0(E)/\mathrm{rk}\,E\leq\mathrm h^0(F)/\mathrm{rk}\,F$ 刻画次数 $d>D$ 的向量丛 $F$ 的半稳定性. GIT 半稳定性判据 (\bookref{9.5.12}) 选取 $M$, 并用对所有子空间 $V'\subseteq V$ 的不等式 $P_{F'}(z)/\dim V'\geq P_F(z)/\dim V$ 刻画商 $[\mathcal O_C^{\oplus P_F(0)}\to F]$ 的 GIT 半稳定性, 其中 $P_F(z)$ 是 Hilbert 多项式, $F':=\im(\mathcal O_C^{\oplus P_F(0)}\to F)$. 定理~\bookref{9.5.1} 的证明来自对这两个判据的仔细联系.
